
\subsubsection{4.1 Dataset}

RedPajama-v2 (Together Computer, 2023), a large-scale web corpus representative of LLM training data.

\begin{table}[htbp]
\centering
\resizebox{\columnwidth}{!}{%
\begin{tabular}{ll}
\toprule
Property & Value \\
\midrule
Source & Common Crawl + curated sources \\
PoC Sample & 5M--10M tokens \\
Language & English \\
Quality Signal & KenLM 5-gram perplexity \\
\bottomrule
\end{tabular}
}
\end{table}

\subsubsection{4.2 Baselines}

\begin{itemize}
\item \textbf{No Filtering (p0):} Raw corpus without perplexity filtering
\item \textbf{RedPajama Defaults:} p30 perplexity threshold, exact deduplication
\item \textbf{CPDR-Optimized:} p50 perplexity threshold, fuzzy\_0.85 deduplication (identified via sweep)
\end{itemize}

\subsubsection{4.3 Implementation Details}

\textbf{Model Architecture:} GPT-2 variants using HuggingFace Transformers.

\textbf{Training Configuration:}
\begin{itemize}
\item Learning rate: 6e-4 with cosine decay
\item Batch size: 512 sequences
\item Training tokens: 5M--10M (PoC validation)
\item Random seeds: 42, 43, 44
\end{itemize}

\textbf{Perplexity Scoring:} KenLM 5-gram model trained on Wikipedia following CCNet methodology.

\subsubsection{4.4 Evaluation Protocol}

Four reasoning benchmarks: HellaSwag, ARC-Easy, PIQA, WinoGrande. Benchmark ensemble score computed as mean accuracy.

\textbf{Note:} Due to lm-evaluation-harness integration issues, mock evaluation was used for PoC validation. Results represent expected patterns based on simulation, not actual benchmark scores.
